\documentclass[11pt]{article}

\usepackage[a4paper,margin=28mm]{geometry}
\usepackage[T1]{fontenc}
\usepackage{lmodern}
\usepackage{microtype}
\usepackage{amsmath,amssymb,amsthm,mathtools}
\usepackage{booktabs}
\usepackage{enumitem}
\usepackage[hidelinks]{hyperref}
\usepackage[nameinlink,noabbrev]{cleveref}

\newtheorem{theorem}{Theorem}[section]
\newtheorem{proposition}[theorem]{Proposition}
\newtheorem{lemma}[theorem]{Lemma}
\newtheorem{corollary}[theorem]{Corollary}
\newtheorem{remark}[theorem]{Remark}
\newtheorem{question}[theorem]{Question}

\crefname{question}{Question}{Questions}
\Crefname{question}{Question}{Questions}
\crefalias{question}{question}
\crefname{lemma}{Lemma}{Lemmas}
\Crefname{lemma}{Lemma}{Lemmas}
\crefname{theorem}{Theorem}{Theorems}
\Crefname{theorem}{Theorem}{Theorems}
\crefname{proposition}{Proposition}{Propositions}
\Crefname{proposition}{Proposition}{Propositions}
\crefname{corollary}{Corollary}{Corollaries}
\Crefname{corollary}{Corollary}{Corollaries}
\crefname{remark}{Remark}{Remarks}
\Crefname{remark}{Remark}{Remarks}

\newcommand{\FBL}{\operatorname{FBL}}
\newcommand{\Id}{\operatorname{Id}}
\newcommand{\pf}{\operatorname{pf}}
\newcommand{\supp}{\operatorname{supp}}
\newcommand{\Homlat}{\operatorname{Hom}_{\mathrm{lat}}}
\newcommand{\R}{\mathbb{R}}
\newcommand{\N}{\mathbb{N}}
\newcommand{\norm}[1]{\left\lVert #1\right\rVert}
\newcommand{\abs}[1]{\left|#1\right|}
\newcommand{\angles}[2]{\left\langle #1,#2\right\rangle}

\title{Positive Factorization Through Free Banach Lattices\\
\large Verification, Literature Status, and Further Reductions}
\author{Research report prepared from the supplied partial solution}
\date{21 June 2026}

\begin{document}
\maketitle

\begin{abstract}
Let $X$ be a Banach lattice, let $\delta_X:X\to\FBL[X]$ be the canonical
linear isometry, and let $\beta_X:\FBL[X]\to X$ be the canonical contractive
lattice homomorphism satisfying $\beta_X\delta_X=\Id_X$.  Oikhberg asked
whether $\beta_X$ always has a positive bounded linear right inverse.

The supplied argument is correct: it gives an explicit positive section for
$\ell_1(I)$ and transfers it to every positive retract of $\ell_1(I)$.  We
also identify the exact scope of that hypothesis: such retracts are precisely,
up to lattice isomorphism, the spaces $\ell_1(A)$.  A literature search through
21 June 2026 found no published or preprint proof or counterexample for the
general question.  The report records several additional affirmative results
and reductions.  In particular, Lotz's positive lifting theorem implies the
answer is affirmative for every AM-space, hence for every $C_0(K)$ and every
$c_0(\Gamma)$.  We prove an equivalent split-quotient formulation, construct a
canonical positive contractive section into $\FBL[X]^{**}$ for every $X$, and
show that the remaining obstruction is precisely weak-star continuity of a
positive dual retraction.  Finally, we analyze the natural construction for
nonatomic $L_1$ and show concretely why it leaves the free Banach lattice and
lands only in its bidual.  The arbitrary Banach-lattice problem remains open.
\end{abstract}

\section{The problem and the conclusion of this investigation}

For a Banach space $E$, the free Banach lattice $\FBL[E]$ comes with a
canonical linear isometry $\delta_E:E\to\FBL[E]$.  Every bounded operator
$T:E\to Y$ into a Banach lattice has a unique lattice-homomorphic extension
\[
 \widehat T:\FBL[E]\longrightarrow Y,
 \qquad \widehat T\delta_E=T,
 \qquad \norm{\widehat T}=\norm T.
\]
If $X$ itself is a Banach lattice, applying this property to $\Id_X$ gives a
contractive lattice homomorphism
\[
 \beta_X:\FBL[X]\longrightarrow X,
 \qquad \beta_X\delta_X=\Id_X.
\]
Thus $\delta_X$ is always a bounded linear section of $\beta_X$, but it is
usually not positive.  The question raised in \cite[Section~3]{Oikhberg2024}
is:

\begin{question}[Oikhberg]\label{q:main}
Does every Banach lattice $X$ admit a positive bounded linear map
$U:X\to\FBL[X]$ such that $\beta_XU=\Id_X$?
\end{question}

Our literature search included the exact wording of the question, the phrases
``positive factorization through $\FBL[X]$'', ``positive right inverse'', and
``split lattice quotient'', citations to \cite{Oikhberg2024}, and recent work
on free, projective, semiprojective, and pre-ordered Banach lattices.  We found
no paper resolving Question~\ref{q:main}.  In particular, the recent construction of
free Banach lattices over pre-ordered Banach spaces in \cite{deJeuJiang2026}
uses a different universal object whose generating map is positive by
definition; it does not produce a positive section into the ordinary
$\FBL[X]$.  Recent work on complemented subspaces and free products likewise
does not settle the question \cite{deHeviaTradacete2025,MartinezTradacete2026}.
A literature search cannot exclude unpublished work, but the best supported
status as of the date of this report is that the general problem remains open.

The main concrete conclusions below are:
\begin{enumerate}[label=(\roman*),leftmargin=2.2em]
\item the supplied partial proof is valid;
\item its positive-$\ell_1$-retract hypothesis is equivalent to being lattice
isomorphic to an atomic $L_1$-space $\ell_1(A)$;
\item all AM-spaces have the desired positive factorization, by a classical
theorem of Lotz;
\item the problem is equivalent to asking whether every linearly split lattice
quotient onto $X$ splits positively;
\item every $X$ has a canonical positive contractive section into
$\FBL[X]^{**}$, and the exact remaining issue is whether a suitable dual
retraction can be chosen weak-star continuous.
\end{enumerate}

\section{Verification of the supplied partial result}

We first reproduce the argument in a slightly sharpened form.  The proof is
correct, including for arbitrary index sets.

\begin{lemma}[The atomic $\ell_1$ lifting]\label{lem:l1}
For every set $I$ there is a positive contraction
\[
 S_I:\ell_1(I)\longrightarrow \FBL[\ell_1(I)]
\]
such that $\beta_{\ell_1(I)}S_I=\Id_{\ell_1(I)}$.
\end{lemma}

\begin{proof}
Let $(e_i)_{i\in I}$ be the canonical atoms of $\ell_1(I)$.  For
$a=(a_i)_{i\in I}\in\ell_1(I)$ put
\[
 S_Ia=\sum_{i\in I}a_i\,\delta_{\ell_1(I)}(e_i)^+.
\]
The net of finite partial sums converges in norm, since
\[
 \sum_{i\in I}|a_i|\,
 \norm{\delta_{\ell_1(I)}(e_i)^+}
 \leq \sum_{i\in I}|a_i|<\infty.
\]
Equivalently, every element of $\ell_1(I)$ has countable support and the
resulting series is absolutely convergent.  Hence $S_I$ is a contraction.  It
is positive because $a_i\geq0$ for all $i$ implies $S_Ia\geq0$.  Finally,
$\beta_{\ell_1(I)}$ is a lattice homomorphism, so
\[
 \beta_{\ell_1(I)}\bigl(\delta(e_i)^+\bigr)
 =\beta_{\ell_1(I)}(\delta(e_i))^+=e_i.
\]
Linearity and norm convergence now give
$\beta_{\ell_1(I)}S_Ia=a$.
\end{proof}

\begin{theorem}[The supplied positive-retract theorem]\label{thm:retract}
Suppose there are positive bounded maps
\[
 J:X\longrightarrow\ell_1(I),
 \qquad Q:\ell_1(I)\longrightarrow X,
 \qquad QJ=\Id_X.
\]
Then there is a positive map $U:X\to\FBL[X]$ satisfying
$\beta_XU=\Id_X$ and
\[
 \norm U\leq \norm Q\,\norm J.
\]
\end{theorem}

\begin{proof}
Let $S_I$ be as in \cref{lem:l1} and let
$\FBL[Q]:\FBL[\ell_1(I)]\to\FBL[X]$ be the lattice homomorphism induced by
the bounded linear map $Q$.  Define
\[
 U=\FBL[Q]S_IJ.
\]
This is positive and has the stated norm bound.  If $Jx=(a_i)$, then
\begin{align*}
 \beta_XUx
 &=\sum_i a_i\,\beta_X\FBL[Q]\bigl(\delta(e_i)^+\bigr)\\
 &=\sum_i a_i\,\beta_X\bigl(\delta_X(Qe_i)^+\bigr)\\
 &=\sum_i a_i Qe_i
 =QJx=x.
\end{align*}
The third equality uses $Qe_i\geq0$.  All series are absolutely convergent in
norm.
\end{proof}

\begin{remark}[Two points worth recording]\label{rem:functoriality}
First, the last calculation is genuinely needed.  The identity
\[
 \beta_X\FBL[T]=T\beta_Y
\]
is automatic when $T:Y\to X$ is a lattice homomorphism, because both sides
are then lattice homomorphisms agreeing on the free generators.  It need not
hold for a merely positive $T$.  Thus the positivity of each $Qe_i$ is the key
special feature of the atomic proof.

Second, the finite-dimensional corollary is correct but was already known in
a stronger form: finite-dimensional Banach lattices have the lattice-lifting
property, in which the section itself can be chosen to be a lattice
homomorphism \cite[Proposition~2.2]{AvilesEtAl2022}.
\end{remark}

\section{What the $\ell_1$-retract hypothesis really means}

The partial theorem appears at first to cover a broad class of positive
retracts.  In fact, that class is exactly the atomic $L_1$ class.

\begin{theorem}[Classification of positive $\ell_1$ retracts]\label{thm:classify}
For a Banach lattice $X$, the following are equivalent.
\begin{enumerate}[label=(\alph*),leftmargin=2.2em]
\item There are a set $I$ and positive bounded maps
$J:X\to\ell_1(I)$ and $Q:\ell_1(I)\to X$ with $QJ=\Id_X$.
\item $X$ is lattice isomorphic to $\ell_1(A)$ for some set $A$.
\end{enumerate}
\end{theorem}

\begin{proof}
Only (a)$\Rightarrow$(b) needs proof.  Let
$\mathbf 1=(1)_{i\in I}\in\ell_\infty(I)$ and define the positive functional
\[
 \varphi=J^*\mathbf 1\in X^*_+.
\]
Set
\[
 \lVert\!\lvert x\rvert\!\rVert=\varphi(|x|)=\norm{J|x|}_{\ell_1(I)}.
\]
This is an equivalent lattice norm.  Indeed,
\[
 \lVert\!\lvert x\rvert\!\rVert
 \leq \norm J\,\norm x,
\]
while positivity of $J$ gives $|Jx|\leq J|x|$, and hence
\[
 \norm x=\norm{QJx}
 \leq \norm Q\,\norm{Jx}_1
 \leq \norm Q\,\lVert\!\lvert x\rvert\!\rVert.
\]
For $x,y\geq0$,
\[
 \lVert\!\lvert x+y\rvert\!\rVert
 =\varphi(x+y)
 =\lVert\!\lvert x\rvert\!\rVert+
  \lVert\!\lvert y\rvert\!\rVert.
\]
Thus $X$, with this equivalent lattice norm, is an abstract AL-space.  By the
Kakutani representation theorem it is lattice isometric to an $L_1(\mu)$
space.

On the other hand, $J$ is an isomorphic embedding of the underlying Banach
space into $\ell_1(I)$, so $X$ has the Schur property.  The Schur property is
unchanged by equivalent norms.  An $L_1(\mu)$ space has the Schur property
only when its measure algebra is purely atomic: a nonzero nonatomic finite
part supports a norm-one weakly null Rademacher sequence.  A purely atomic
$L_1$ space is lattice isometric, after absorbing the atom weights, to
$\ell_1(A)$ for a suitable set $A$.  This proves (b).  The converse follows by
taking $I=A$ and the identity maps.
\end{proof}

\begin{corollary}
The new content of \cref{thm:retract}, up to lattice isomorphism, is precisely
the positive factorization for the spaces $\ell_1(A)$.
\end{corollary}

This also corrects one limitation stated in the supplied packet.  The packet
said that it gave no answer for $c_0$.  Although $c_0$ is not in the class of
\cref{thm:classify}, it does have the desired factorization, and even the
stronger lattice-lifting property; see the next section.

\section{Affirmative classes already available in the literature}

\subsection{The stronger lattice-lifting property}

Avil\'es, Mart\'inez-Cervantes, Rodr\'iguez, and Tradacete call $X$ a lattice
lifting space if there is a lattice homomorphism
$A:X\to\FBL[X]$ with $\beta_XA=\Id_X$ \cite{AvilesEtAl2022}.  This is stronger
than Question~\ref{q:main}.  Their results include the following.
\begin{itemize}[leftmargin=2em]
\item Every projective Banach lattice has the lattice-lifting property; in
particular this includes finite-dimensional Banach lattices, $\ell_1$, and
$C[0,1]$.
\item Every free Banach lattice $\FBL[E]$ has the isometric lattice-lifting
property.
\item Every Banach space with a $1$-unconditional Schauder basis, equipped
with its coordinate order, has the isometric lattice-lifting property.
Consequently, the classical separable spaces $c_0$ and $\ell_p$ have it.
\end{itemize}
The same paper shows that the stronger property fails for nonatomic $L_p$
spaces and for $c_0(\Gamma)$ and $\ell_p(\Gamma)$ when $\Gamma$ is uncountable.
Those negative results do not answer Question~\ref{q:main}, since a positive section
need not preserve disjointness.

\subsection{A classical theorem of Lotz gives all AM-spaces}

The following consequence of a 1975 theorem of Lotz appears not to have been
noted in the discussion surrounding Question~\ref{q:main}.

\begin{theorem}[Lotz]\label{thm:lotz}
Every AM-space $X$ admits a positive bounded map
$U:X\to\FBL[X]$ such that $\beta_XU=\Id_X$.
Consequently the same is true for every Banach lattice lattice-isomorphic to
an AM-space.
\end{theorem}

\begin{proof}
Lotz proved that every AM-space is \emph{$d$-projective}: if
$q:Y\to Z$ is a nearly interval-preserving metric surjection of Banach
lattices and $T:X\to Z$ is positive, then $T$ has a positive bounded lifting
through $q$ \cite[Proposition~3.3]{Lotz1975}.

We verify the hypotheses for $q=\beta_X$.  It is a metric surjection: for
$x\in X$,
\[
 \inf\{\norm f:\beta_Xf=x\}=\norm x,
\]
because $\beta_X$ is contractive and $\delta_Xx$ is a preimage of norm
$\norm x$.  It is in fact interval preserving.  If $f\geq0$ in $\FBL[X]$ and
$0\leq x\leq\beta_Xf$, choose $g\in\FBL[X]$ with $\beta_Xg=x$ and put
$h=f\wedge g^+$.  Then $0\leq h\leq f$ and, since $\beta_X$ is a lattice
homomorphism,
\[
 \beta_Xh=\beta_Xf\wedge(\beta_Xg)^+
 =\beta_Xf\wedge x=x.
\]
Apply Lotz's theorem to the positive map $T=\Id_X:X\to X$.
\end{proof}

\begin{corollary}\label{cor:am}
The answer to Question~\ref{q:main} is affirmative for every $C_0(K)$ space and every
$c_0(\Gamma)$, with arbitrary locally compact $K$ and arbitrary set $\Gamma$.
In particular, uncountable $c_0(\Gamma)$ has a positive section although it
has no lattice-homomorphic section.
\end{corollary}

For orientation, the present state of several standard examples is summarized
below.  ``Lattice section'' means the stronger property from
\cite{AvilesEtAl2022}.

\begin{center}
\begin{tabular}{@{}lll@{}}
\toprule
Space or class & Positive section & Lattice section \\
\midrule
$\ell_1(\Gamma)$ & yes, explicit & no in general if $\Gamma$ uncountable \\
$c_0(\Gamma)$ & yes, Lotz & no in general if $\Gamma$ uncountable \\
$C_0(K)$ / AM-spaces & yes, Lotz & not in general \\
$1$-unconditional Schauder basis & yes & yes \\
$\FBL[E]$ & yes & yes, isometric \\
finite-dimensional lattices & yes & yes \\
nonatomic $L_1$ & open & no \\
\bottomrule
\end{tabular}
\end{center}

\section{An exact split-quotient reformulation}

The next proposition turns the free-lattice question into a problem purely
about split lattice quotients.  It is the positive analogue of
\cite[Proposition~2.1]{AvilesEtAl2022}.

Define
\[
 \pf(X)=\inf\bigl\{\norm U:
 U:X\to\FBL[X]\text{ is positive and }\beta_XU=\Id_X\bigr\},
\]
with $\pf(X)=\infty$ if no such map exists.

\begin{proposition}[Split-quotient characterization]\label{prop:split}
Let $X$ be a Banach lattice and $\lambda\geq1$.  The following are equivalent.
\begin{enumerate}[label=(\alph*),leftmargin=2.2em]
\item $\pf(X)\leq\lambda$.
\item For every $\varepsilon>0$, whenever $Y$ is a Banach lattice,
$q:Y\to X$ is a surjective lattice homomorphism, and $T:X\to Y$ is a
bounded linear map with $qT=\Id_X$, there is a positive map $S:X\to Y$
such that
\[
 qS=\Id_X,
 \qquad \norm S\leq(\lambda+\varepsilon)\norm T.
\]
\end{enumerate}
\end{proposition}

\begin{proof}
Assume (a), fix $\varepsilon>0$, and choose a positive section
$U:X\to\FBL[X]$ with $\norm U\leq\lambda+\varepsilon$.  By the universal
property, $T$ extends to a lattice homomorphism
$\widehat T:\FBL[X]\to Y$ with $\norm{\widehat T}=\norm T$.  Both
$q\widehat T$ and $\beta_X$ are lattice homomorphisms from $\FBL[X]$ to $X$,
and
\[
 q\widehat T\delta_X=qT=\Id_X=\beta_X\delta_X.
\]
Uniqueness gives $q\widehat T=\beta_X$.  Hence $S=\widehat TU$ is positive,
$qS=\Id_X$, and $\norm S\leq(\lambda+\varepsilon)\norm T$.

Conversely, apply (b), for every $\varepsilon>0$, with
$Y=\FBL[X]$, $q=\beta_X$, and $T=\delta_X$.  This gives
$\pf(X)\leq\lambda$.
\end{proof}

\begin{corollary}[Ideal formulation]\label{cor:ideal}
A positive section for $\beta_X$ exists if and only if every lattice quotient
$q:Y\to X$ whose kernel is complemented as a Banach subspace has a positive
right inverse.  Equivalently, every such kernel is the kernel of a positive
projection on $Y$.
\end{corollary}

\begin{proof}
A quotient has a bounded linear right inverse exactly when its kernel is
linearly complemented.  A positive section $S$ produces the positive
projection $P=Sq$, with $\ker P=\ker q$.  Conversely, if $P$ is positive and
$\ker P=\ker q$, then $q|_{P(Y)}$ is bijective and its inverse is positive:
for $x\geq0$, choose $y\geq0$ with $qy=x$ and observe that
$(q|_{P(Y)})^{-1}x=Py\geq0$.
\end{proof}

Thus a counterexample to Question~\ref{q:main} can be sought without mentioning free
Banach lattices: it is enough, and necessary, to find a linearly split lattice
quotient with no positive splitting.  We found no such example in the
literature.  Conversely, a theorem asserting that every linearly split lattice
quotient splits positively would solve the problem affirmatively.

\section{A universal positive lifting into the bidual}

There is always a positive section after passing to the bidual.  This is both a
partial affirmative result and a precise description of the obstruction.

\begin{theorem}[Canonical bidual section]\label{thm:bidual}
For every Banach lattice $X$ there is a positive contraction
\[
 U_0:X\longrightarrow\FBL[X]^{**}
\]
such that
\[
 \beta_X^{**}U_0=J_X,
\]
where $J_X:X\to X^{**}$ is the canonical embedding.
\end{theorem}

\begin{proof}
Put $F=\FBL[X]$ and $I=\ker\beta_X$.  Since $F/I$ is lattice isometric to
$X$, the adjoint
\[
 \beta_X^*:X^*\longrightarrow F^*
\]
identifies $X^*$ lattice isometrically with the annihilator $I^\perp$.  The
annihilator of a closed ideal is a band in the dual Banach lattice $F^*$.
Dual Banach lattices are Dedekind complete, so there is a band projection
\[
 P_I:F^*\longrightarrow I^\perp.
\]
It is positive and contractive.  Define
\[
 R_0=(\beta_X^*)^{-1}P_I:F^*\longrightarrow X^*.
\]
Then $R_0$ is a positive contraction and $R_0\beta_X^*=\Id_{X^*}$.  Taking
adjoints yields
\[
 \beta_X^{**}R_0^*=\Id_{X^{**}}.
\]
The map $U_0=R_0^*J_X$ has all the required properties.
\end{proof}

\begin{proposition}[The exact weak-star criterion]\label{prop:wstar}
There is a positive section $U:X\to\FBL[X]$ for $\beta_X$ if and only if there
is a positive weak-star-to-weak-star continuous map
\[
 R:\FBL[X]^*\longrightarrow X^*
\]
satisfying $R\beta_X^*=\Id_{X^*}$.  In that case one can take $R=U^*$.
\end{proposition}

\begin{proof}
If $U$ exists, $U^*$ is positive, weak-star continuous, and
$U^*\beta_X^*=\Id$.  Conversely, a bounded weak-star-to-weak-star continuous
operator $R$ has a preadjoint $U:X\to\FBL[X]$.  Positivity of $R$ implies
positivity of $U$, and taking adjoints in $R\beta_X^*=\Id$ gives
$\beta_XU=\Id$.
\end{proof}

The band retraction $R_0$ in \cref{thm:bidual} always exists.  The unresolved
point is that a band projection on a dual Banach lattice need not be weak-star
continuous.  This is not merely a technicality; the canonical $R_0$ behaves
very badly on nonatomic examples.

For $x^*\in X^*$, let $a_{x^*}\in\FBL[X]^*$ be evaluation at $x^*$ in the
standard functional representation:
\[
 a_{x^*}(f)=f(x^*).
\]
When $x^*\neq0$, this is an atom of $\FBL[X]^*$.

\begin{proposition}[Action of the canonical band projection on atoms]
\label{prop:atoms}
Let $P_I$ be the band projection used in \cref{thm:bidual}.  Then
\[
 P_Ia_{x^*}=
 \begin{cases}
 a_{x^*}=\beta_X^*x^*,&\text{if }x^*:X\to\R\text{ is a lattice homomorphism},\\
 0,&\text{otherwise}.
 \end{cases}
\]
Consequently,
\[
 U_0(x)(a_{x^*})=
 \begin{cases}
 x^*(x),&x^*\in\Homlat(X,\R),\\
 0,&x^*\notin\Homlat(X,\R).
 \end{cases}
\]
\end{proposition}

\begin{proof}
A band projection sends an atom either to itself or to zero.  We have
$P_Ia_{x^*}=a_{x^*}$ exactly when $a_{x^*}\in I^\perp=\beta_X^*(X^*)$.
If $a_{x^*}=\beta_X^*\phi$, evaluation on $\delta_X(X)$ gives $\phi=x^*$.
Thus membership is equivalent to
$a_{x^*}=\beta_X^*x^*$.  If $x^*$ is a lattice homomorphism, both sides are
lattice homomorphisms on $\FBL[X]$ extending the same map on $X$, so uniqueness
of the free extension gives equality.  Conversely, if equality holds, then
$x^*\beta_X$ is a lattice homomorphism; surjectivity of $\beta_X$ implies that
$x^*$ itself is a lattice homomorphism.  The formula for $U_0$ follows from
the definitions.
\end{proof}

\begin{corollary}\label{cor:canonical-fails}
If the real-valued lattice homomorphisms on $X$ do not separate points, then
for every nonzero $x$ annihilated by all of them, the element $U_0x$ does not
belong to the canonical copy of $\FBL[X]$ inside $\FBL[X]^{**}$.  In
particular, this happens for every nonzero $x\in L_1[0,1]$.
\end{corollary}

\begin{proof}
If $U_0x=J_Ff$ for $f\in F=\FBL[X]$, then \cref{prop:atoms} gives
$f(x^*)=0$ for every $x^*\in X^*$.  The functional representation of
$\FBL[X]$ is faithful, hence $f=0$.  But then
$J_Xx=\beta_X^{**}U_0x=0$, a contradiction.
\end{proof}

This shows why the automatic bidual lifting does not settle the original
problem, especially on nonatomic lattices.

\section{Stability properties}

The following elementary permanence results enlarge any affirmative class.

\begin{proposition}\label{prop:stability}
The positive factorization property has the following stability properties.
\begin{enumerate}[label=(\alph*),leftmargin=2.2em]
\item It is invariant under lattice isomorphism.
\item Suppose $Y$ has the property, $J:X\to Y$ is positive,
$Q:Y\to X$ is a lattice homomorphism, and $QJ=\Id_X$.  Then $X$ has the
property.
\item If each $X_i$ has a positive section $U_i$ with
$\sup_i\norm{U_i}<\infty$, then the $\ell_1$-sum
$X=(\bigoplus_i X_i)_1$ has the property.
\end{enumerate}
\end{proposition}

\begin{proof}
For (a), if $A:X\to Y$ is a lattice isomorphism and $U_Y$ is a positive
section, then
\[
 U_X=\FBL[A^{-1}]U_YA
\]
is a positive section; naturality of $\beta$ for lattice homomorphisms gives
the identity.

For (b), put $U_X=\FBL[Q]U_YJ$.  Since $Q$ is a lattice homomorphism,
\[
 \beta_X\FBL[Q]=Q\beta_Y,
\]
and therefore $\beta_XU_X=Q\beta_YU_YJ=QJ=\Id_X$.

For (c), let $j_i:X_i\to X$ be the coordinate lattice embeddings and set
\[
 U(x_i)_i=\sum_i \FBL[j_i]U_ix_i.
\]
The series converges absolutely, with
$\norm{U(x_i)}\leq (\sup_i\norm{U_i})\sum_i\norm{x_i}$.  It is positive, and
naturality gives
\[
 \beta_XU(x_i)=\sum_i j_i\beta_{X_i}U_ix_i=(x_i).
\]
\end{proof}

Part (b) explains why lattice retracts inherit the property.  The supplied
\cref{thm:retract} is stronger in one special direction: there $Q$ is only
positive, but the atomic coordinate formula compensates for the lack of
naturality noted in \cref{rem:functoriality}.

\section{Why the obvious nonatomic $L_1$ construction fails}

The atomic formula suggests replacing a sum over atoms by an integral.  This
produces a bounded positive map into the ambient space of homogeneous
functions, but not into the free Banach lattice.

Let $X=L_1[0,1]$, so $X^*=L_\infty[0,1]$.  For $f\in L_1$ define
\[
 \Gamma_f(g)=\int_0^1 f(t)g(t)^+\,dt,
 \qquad g\in L_\infty.
\]
The map $f\mapsto\Gamma_f$ is linear and positive with respect to pointwise
order on homogeneous functions.  Moreover, for any finite family
$g_1,\dots,g_n\in L_\infty$,
\[
 \sum_{j=1}^n|\Gamma_f(g_j)|
 \leq \norm{f}_1\,
 \norm{\sum_{j=1}^n|g_j|}_\infty.
\]
Thus $\Gamma_f$ satisfies exactly the norm estimate one would expect for an
element of $\FBL[L_1]$.

Nevertheless, $\Gamma_1\notin\FBL[L_1]$.  Let $(r_n)$ be the Rademacher
functions.  Then $r_n\to0$ in $\sigma(L_\infty,L_1)$: this is immediate first
against dyadic simple functions and then follows for all $L_1$ functions by
density.  But
\[
 \Gamma_1(r_n)=\int_0^1r_n^+\,dt=\frac12
 \qquad(n\in\N).
\]
Elements of $\FBL[E]$ are weak-star continuous on bounded subsets of $E^*$ in
the standard functional representation.  Hence $\Gamma_1$ lies outside
$\FBL[L_1]$.

There is also a finite-dimensional approximation that makes the obstruction
visible.  For a finite measurable partition
$\mathcal P=\{E_1,\dots,E_m\}$ with $\mu(E_i)>0$, define
\[
 U_{\mathcal P}f
 =\sum_{i=1}^m
 \left(\int_{E_i}f\,d\mu\right)
 \delta_{L_1}\left(\frac{\mathbf1_{E_i}}{\mu(E_i)}\right)^+.
\]
Then $U_{\mathcal P}:L_1\to\FBL[L_1]$ is a positive contraction and
\[
 \beta_{L_1}U_{\mathcal P}f
 =\mathbb E(f\mid\mathcal P).
\]
Along refining dyadic partitions, the right-hand side converges to $f$ in
$L_1$.  However, for $f=1$ and $g\in L_\infty$,
\[
 U_{\mathcal P}1(g)
 =\sum_i\mu(E_i)
 \left(\frac1{\mu(E_i)}\int_{E_i}g\,d\mu\right)^+,
\]
which converges pointwise to $\int g^+=\Gamma_1(g)$.  The bounded net therefore
converges only in a bidual/pointwise sense, to an object outside the free
Banach lattice.  This is a concrete manifestation of the weak-star
continuity obstruction in \cref{prop:wstar}.

This calculation is not a counterexample: it rules out the most natural
integral formula, not every possible positive section for $L_1[0,1]$.

\section{Final assessment and viable routes forward}

The supplied result is mathematically sound.  Its principal proof requires no
repair, only two contextual corrections: the finite-dimensional case was
known more strongly, and $c_0$ is already affirmative.  The exact scope of the
positive-retract hypothesis is the class of atomic $L_1$ spaces
$\ell_1(A)$.

The strongest additional full theorem located is the AM-space consequence of
Lotz's positive lifting theory, which yields all $C_0(K)$ and all
$c_0(\Gamma)$.  The general problem, including the basic nonatomic test case
$L_1[0,1]$, remains unresolved in the literature we found.

The reductions above suggest two focused approaches.
\begin{enumerate}[label=(\arabic*),leftmargin=2.2em]
\item By \cref{prop:split}, construct a Banach lattice $Y$ and a linearly
complemented ideal $I$ such that the quotient map $Y\to Y/I$ has no positive
right inverse.  This would give a counterexample immediately.
\item By \cref{prop:wstar}, seek a positive weak-star continuous retraction of
$\FBL[X]^*$ onto the band $\beta_X^*(X^*)$.  The canonical band projection is
positive and contractive but is generally not weak-star continuous; any
successful proof must replace it by a different positive retraction or prove
additional structure forcing weak-star continuity.
\end{enumerate}

The finite-partition computation for $L_1$ shows that a bare compactness or
local finite-dimensional argument is unlikely to suffice: it naturally
produces the canonical bidual section rather than an element of
$\FBL[L_1]$.

\begin{thebibliography}{99}

\bibitem{ART2018}
A. Avil\'es, J. Rodr\'iguez, and P. Tradacete,
\emph{The free Banach lattice generated by a Banach space},
J. Funct. Anal. 274 (2018), 2955--2977.

\bibitem{AvilesEtAl2022}
A. Avil\'es, G. Mart\'inez-Cervantes, J. Rodr\'iguez, and P. Tradacete,
\emph{A Godefroy--Kalton principle for free Banach lattices},
Israel J. Math. 247 (2022), 433--458;
\url{https://arxiv.org/abs/2011.04639}.

\bibitem{deHeviaTradacete2025}
D. de Hevia and P. Tradacete,
\emph{Complemented subspaces of Banach lattices},
preprint (2025), \url{https://arxiv.org/abs/2505.24084}.

\bibitem{deJeuJiang2026}
M. de Jeu and X. Jiang,
\emph{Free Banach lattices over pre-ordered Banach spaces},
arXiv:2408.06630, version 3 (22 April 2026),
\url{https://arxiv.org/abs/2408.06630}.

\bibitem{Lotz1975}
H. P. Lotz,
\emph{Extensions and liftings of positive linear mappings on Banach lattices},
Trans. Amer. Math. Soc. 211 (1975), 85--100.

\bibitem{MartinezTradacete2026}
G. Mart\'inez-Fern\'andez and P. Tradacete,
\emph{Free products of Banach lattices},
arXiv:2605.28988 (2026),
\url{https://arxiv.org/abs/2605.28988}.

\bibitem{MeyerNieberg}
P. Meyer-Nieberg,
\emph{Banach Lattices}, Springer, Berlin, 1991.

\bibitem{Oikhberg2024}
T. Oikhberg,
\emph{Geometry of unit balls of free Banach lattices, and its applications},
J. Funct. Anal. 286 (2024), no.~8, Article 110351;
\url{https://arxiv.org/abs/2303.05209}.

\bibitem{FBLMemoir}
T. Oikhberg, M. A. Taylor, P. Tradacete, and V. G. Troitsky,
\emph{Free Banach lattices},
preprint, \url{https://arxiv.org/abs/2210.00614}.

\end{thebibliography}

\end{document}
